% !PRINTED 638
\weberpage{638}
%<<CONT
unter $\vartheta$ ein Product aus Potenzen der Zahlen $\varrho$ versteht, die
nach den Formeln (24) aus den in $\varphi$ enthaltenen Primfactoren $\mathfrak{p}$
abgeleitet sind.

\begin{center}
§. 175.

\gesperrt{Die Einheitswurzeln im Körper $\Omega_m$.}
\end{center}

Für die weiter zu machenden Anwendungen ist eine genauere
Kenntniss der Einheiten des Körpers $\Omega_m$ erforderlich, die ja
auch an sich von grossem Interesse ist. Wir beschäftigen uns
hier nicht mit den functionalen, sondern nur mit den numerischen
Einheiten, worunter, wie wir uns erinnern, \gesperrt{ganze Zahlen des
Körpers $\Omega_m$ zu verstehen sind, deren Norm $\pm 1$ ist,
oder eine ganze Zahl, deren reciproker Werth gleichfalls ganz ist.}

Zu den Einheiten gehören sicher alle in $\Omega_m$ enthaltenen
Einheitswurzeln; wir können aber leicht beweisen, dass diese
Einheitswurzeln nur Potenzen von $r$ sein können, mit positivem
und negativem Vorzeichen.

Es sei nämlich $\varrho$ eine in $\Omega_m$ enthaltene Einheitswurzel, und
\weq{1}{\varrho = \varphi(r).}

Ist $q^h$ die höchste im Grade von $\varrho$ aufgehende Potenz von $q$,
so ist $\varrho^{q^h}$ eine Einheitswurzel, deren Grad nicht durch $q$ theilbar
ist, und
\weq{2}{\varrho^{q^h} = [\varphi(r)]^{q^h}.}

Wegen der Irreducibilität der Kreistheilungsgleichung im
weiteren Sinne\footnote{Vgl. Bd.~I, §.\,134 und den Nachtrag zu diesem Bande.} können daher in (2) die sämmtlichen Substitutionen
 $(r, r^n)$ gemacht werden, ohne dass die linke Seite sich
ändert, und folglich ist $\varrho^{q^h}$ eine rationale Zahl, und da es eine
Einheit ist, muss es $= \pm 1$ sein. Mithin ist $\varrho$, vom Vorzeichen
abgesehen, einer Einheitswurzel vom Grade $q^h$ oder, wenn $q = 2$
ist, vom Grade $q^{h+1}$ gleich. Es ist nachzuweisen, dass $q^h$ oder
$q^{h+1}$ nicht grösser als $m$ sein kann. Dies ergiebt sich aber
aus (1). Denn wäre $\varrho$ eine Einheitswurzel höheren als $m^{\text{ten}}$
Grades, so wäre $r$ eine Potenz von $\varrho$, und $\varrho$ müsste einer irreduciblen
 Gleichung höheren Grades als $r$, also auch höheren
Grades als $\varphi(r)$ genügen, was nach (1) ein Widerspruch ist. Also
haben wir den Satz:
%CONT>>
